\documentclass{article}
\usepackage{amsmath,amssymb}
\usepackage{xurl}
\usepackage[colorlinks=true,linkcolor=blue,citecolor=blue,urlcolor=blue]{hyperref}
% Shared commands for notes in docs/paper-gaps/.

\DeclareUnicodeCharacter{2080}{\ifmmode{}_{0}\else\textsubscript{0}\fi}
\DeclareUnicodeCharacter{2081}{\ifmmode{}_{1}\else\textsubscript{1}\fi}
\DeclareUnicodeCharacter{2082}{\ifmmode{}_{2}\else\textsubscript{2}\fi}
\DeclareUnicodeCharacter{2099}{\ifmmode{}_{n}\else\textsubscript{n}\fi}

\newcommand{\Lean}{\textsc{Lean}}
\ExplSyntaxOn
\NewDocumentCommand{\leanid}{m}{
  \ifmmode
    \textsf{\detokenize{#1}}
  \else
    \group_begin:
    \urlstyle{sf}
    \tl_set:Nn \l_tmpa_tl {#1}
    \tl_replace_all:Nnn \l_tmpa_tl {\_} {_}
    \exp_args:NV \path \l_tmpa_tl
    \group_end:
  \fi
}
\ExplSyntaxOff
\providecommand{\tr}{\operatorname{tr}}
\newcommand{\tnleanIssueBase}{https://github.com/LionSR/TNLean/issues/}
\newcommand{\tnleanPrBase}{https://github.com/LionSR/TNLean/pull/}
\newcommand{\tnleanDocBase}{https://sirui-lu.com/TNLean/docs/}
\newcommand{\ghissue}[1]{\href{\tnleanIssueBase#1}{GitHub issue \##1}}
\newcommand{\ghpr}[1]{\href{\tnleanPrBase#1}{PR \##1}}
\newcommand{\doclink}[1]{\href{\tnleanDocBase#1.html}{\path{#1}}}

% Dirac notation and the identity operator, as in blueprint/src/macros/common.tex.
% \providecommand keeps note-local definitions authoritative.
\usepackage{dsfont}
\providecommand{\ket}[1]{|#1\rangle}
\providecommand{\bra}[1]{\langle #1|}
\providecommand{\braket}[2]{\langle #1 | #2 \rangle}
\providecommand{\ketbra}[2]{|#1\rangle\!\langle #2|}
\providecommand{\Id}{\mathds{1}}
\providecommand{\one}{\mathds{1}}
\providecommand{\C}{\mathbb{C}}
\providecommand{\MN}[1]{M_{#1}(\C)}
\providecommand{\MD}{M_D(\C)}

% Verdict marker: \gapnote{<kind>}{<status>}. Kind is the policy
% classification (clarification, local-correction, scope-restriction,
% unfaithful, false-source, open-gap); status is open, wip (elimination
% underway), resolved, or historical. Machine-read by the site index;
% prints a verdict line.
\newcommand{\gapnote}[2]{\par\noindent\textbf{Verdict:} #1 (#2).\par}

\setlength{\emergencystretch}{2em}
\title{The Threshold of the Spectral Cut-offs in Hastings' Approximate Ground-State Projector}
\author{}
\date{October 22, 2026}
\begin{document}
\maketitle
\gapnote{local-correction}{resolved}

\section{The source statement}

In the appendix of~\cite{Hastings2007AreaLaw}, the localized filtered terms
$M_L$ and $M_R$ satisfy $\lVert M_L\Psi_0\rVert,\lVert M_R\Psi_0\rVert\le
    \delta$ with $\delta=\mathcal O(\Delta E^{-1}J^2e^{-l/3\xi})$, and $O_L$
($O_R$) is defined as the projection onto the eigenvectors of $M_L$ ($M_R$)
with eigenvalue less than or equal to
$\theta=\Delta E^{-1}J^2e^{-l/6\xi}$.\footnote{\path{References/0705.2024/law16.tex},
    lines~829--841.}
The source then uses, in the estimate leading to its final bound,
\begin{align}
    \lVert(M_L+M_R)O_LO_R\rVert
     & \le 2\theta.
    \label{eq:hastings07_spectral_cutoff_source_use}
\end{align}

\section{The point at issue}

The terms $M_L$ and $M_R$ are shifted to have ground-state expectation zero,
so they are not positive in general. If $M_L$ has an eigenvalue $\lambda<
    -\theta$, its eigenvector lies in the range of the printed projection $O_L$,
and $\lVert M_LO_L\rVert\ge\lvert\lambda\rvert>\theta$. The bound
(\ref{eq:hastings07_spectral_cutoff_source_use}) then fails.

\section{The corrected statement}

Let $O_\theta(M)$ be the projection onto the eigenvectors of $M$ with
eigenvalue of absolute value at most $\theta$. For Hermitian $M$, $\theta>0$
and every vector $\Psi$,
\begin{align}
    \lVert\Psi-O_\theta(M)\Psi\rVert
     & \le\frac{\lVert M\Psi\rVert}{\theta},
    \label{eq:hastings07_spectral_cutoff_obd}
    \\
    \lVert M\,O_\theta(M)\rVert
     & \le\theta.
    \label{eq:hastings07_spectral_cutoff_norm}
\end{align}
The first bound is the source's estimate on the moved ground vector, with
constant one; the second gives
(\ref{eq:hastings07_spectral_cutoff_source_use}) once $O_L$ and $O_R$
commute, which holds because they are supported on disjoint
regions.\footnote{\leanid{TNLean.MPS.AreaLaw.Hastings.norm_sub_spectralCutoff_le}
    and \leanid{TNLean.MPS.AreaLaw.Hastings.norm_mul_spectralCutoff_le} in
    \doclink{TNLean/MPS/AreaLaw/Hastings/SpectralCutoff}.}

\section{Conclusion}

Replacing the one-sided threshold by the absolute threshold changes no
constant and no conclusion of the source; it makes the estimate it uses hold
for terms that are not positive.

\bibliographystyle{alpha}
\bibliography{references}

\end{document}
